% Generated draft: requires mathematical and source review.
\documentclass[11pt]{article}
\usepackage[margin=25mm]{geometry}
\usepackage{fontspec}
\setmainfont{Latin Modern Roman}
\usepackage{amsmath,amssymb,mathtools,unicode-math}
\setmathfont{Latin Modern Math}
\usepackage{enumitem,needspace,xurl,hyperref}
\setlength{\parindent}{0pt}
\setlength{\parskip}{5pt}
\setlength{\emergencystretch}{3em}
\newcommand{\N}{\mathbb{N}}
\newcommand{\Z}{\mathbb{Z}}
\newcommand{\Q}{\mathbb{Q}}
\newcommand{\R}{\mathbb{R}}
\newcommand{\C}{\mathbb{C}}
\newcommand{\F}{\mathbb{F}}
\newcommand{\A}{\mathbb{A}}
\newcommand{\PP}{\mathbb{P}}
\newcommand{\E}{\mathbb{E}}
\newcommand{\eps}{\varepsilon}
\newcommand{\dd}{\,\mathrm{d}}
\begin{document}
\section*{UnsolvedMath ID: 30001805}
\textbf{Source problem code:} catalogue entry
\subsection{Definitions and mathematical statement}
A hyperbolic metric is a Riemannian metric of constant Gaussian curvature \(-1\). The source asks three related questions about surfaces, but does not specify a single conjecture or a complete set of hypotheses.

For a fixed surface \(S\), let \(\mathcal H(S)\) denote a class of admissible hyperbolic structures, once the geometric conditions and equivalence relation have been specified. For \(X\in\mathcal H(S)\), the usual closed-geodesic systole is
\[
\operatorname{sys}(X)=\inf\{\ell_X(\gamma):\gamma\text{ is a nonconstant closed geodesic in }X\},
\]
provided the indicated collection is nonempty. Here \(\ell_X(\gamma)\) is its hyperbolic length. If surfaces with boundary are allowed, whether boundary geodesics enter this definition must also be specified. Completeness, finite area, punctures, boundary conditions, and restrictions on boundary lengths are not supplied by the record.

\begin{enumerate}
\item Describe the geometry of surfaces as their genus \(g\) becomes large, and describe the geometry of a typical surface of large genus. The source specifies neither a family of geometric structures nor a probability measure defining ``typical,'' nor the geometric quantities or mode of convergence to be studied. Consequently, no definite probabilistic or asymptotic assertion can be extracted.
\item For a fixed topological type \(S\), determine which admissible hyperbolic structures have the largest possible systole. After specifying \(\mathcal H(S)\) and the systole convention, the extremal question is to determine
\[
M(S;\mathcal H)=\sup_{X\in\mathcal H(S)}\operatorname{sys}(X)
\]
and identify all \(X_*\in\mathcal H(S)\) satisfying
\[
\operatorname{sys}(X_*)=M(S;\mathcal H).
\]
This expresses the question without asserting that the supremum is finite or attained. The source does not restrict the topological type to closed surfaces or supply the admissibility conditions needed to make this optimization problem definite.
\item Determine what the geometry of hyperbolic surfaces implies about metrics on Teichm\"uller space. Once the surface type and geometric conditions are fixed, Teichm\"uller space \(\mathcal T(S)\) is understood as a space of marked structures: a marking identifies the underlying topological surface with \(S\), and two marked structures are equivalent when an isometry between them respects the markings up to isotopy. The orientation and boundary conventions must be fixed where applicable. The source does not identify a metric on \(\mathcal T(S)\), a class of metrics, or a proposed relation or estimate to establish.
\end{enumerate}
These are source-level research questions. The unresolved choices above prevent their replacement by a uniquely specified theorem or conjecture.
\subsection{Short English statement}
The record asks how large-genus surfaces behave, which hyperbolic structures of a fixed topological type maximize the systole, and how surface geometry informs metrics on Teichm\"uller space. A precise formulation still requires a class of surfaces, geometric boundary conditions, a meaning of ``typical,'' and a specified metric or target relation on Teichm\"uller space.
\subsection{Sources}
\begin{itemize}
\item Supplied local registry record: permanent UnsolvedMath ID \(30001805\), title ``Typical Geometry and Systoles of Large-Genus Surfaces,'' source problem code \texttt{catalogue entry}, in the Oberwolfach Reports Open Problems collection. The supplied proposed year is 2011; no individual proposer is identified.
\item Supplied authoritative report citation: \emph{Arithmetic Groups vs. Mapping Class Groups: Similarities, Analogies and Differences} (2011), Oberwolfach Reports, pp.\ 1637--1708, DOI \texttt{10.4171/owr/2011/30}. Source URL and supplied literature-source link, labeled ``original Oberwolfach report'': \url{https://doi.org/10.4171/owr/2011/30}.
\item Supplied literature assessment dated \texttt{2026-08-21}: the record labels the problem ``open'' and describes a broad question with substantial systolic and moduli-space results but no settled general answer to the extracted formulation. This is supplied dated curation, not an independently verified assessment; the source requests independent verification of time-sensitive status and cited claims.
\item The supplied statement status is \texttt{corrected\_verified}. Its verification note reports that the mathematical content was checked against the official Oberwolfach report and that notation, context, and question form were normalized without changing the source-level claim. This verification is reported from the registry and has not been independently performed for this draft.
\end{itemize}
\end{document}
